\section{Gibbs states of nearest-neighbor commuting Hamiltonians}

\begin{definition}[GSNNCH]\label{def:mpdo_source_gsnnch}
    \lean{MPOTensor.GSNNCHData}
    \lean{MPOTensor.HasGSNNCHFormAt}
    \lean{MPOTensor.HasGSNNCHForm}
    \lean{MPOTensor.IsGSNNCHAt}
    \lean{MPOTensor.IsGSNNCH}
    \leanok
    \uses{def:mpo_tensor}
    A density operator on a periodic chain is a Gibbs state of a
    nearest-neighbor commuting Hamiltonian if
    \[
        \rho^{(N)}
        \propto
        \bigoplus_x n_x
        \exp\!\left(-\sum_{j=1}^N\tau_j(h^{(x)})\right),
        \qquad n_x\in\N,
    \]
    where the one-site spaces belonging to distinct labels $x$ are
    orthogonal and
    $[h^{(x)},\tau_1(h^{(x)})]=0$. With the projector-limit convention of
    \cite[Definition~4.8, lines~829--850]{Cirac2016MPDO_arXiv}, this is
    equivalently
    \[
        \rho^{(N)}
        \propto
        \bigoplus_x n_x
        \prod_{j=1}^N\tau_j(B^{(x)}),
        \qquad B^{(x)}\geq0,
    \]
    with commuting neighboring translates. The orthogonal sum and the
    multiplicities are part of the definition.
\end{definition}

The following auxiliary single-bond presentation does not make the sector
labels or multiplicities explicit. Its bond acts on the full two-site space.
It gives the one-sector case of
Definition~\ref{def:mpdo_source_gsnnch}; the converse comparison is not used.

\begin{definition}[Single-bond commuting-form data for an MPDO]
    \label{def:mpdo_commuting_form_data}
    \lean{MPOTensor.CommutingFormData}
    \leanok
    \uses{def:mpo_tensor}
    For a fixed chain length $N \ge 2$, this consists of one positive
    semidefinite two-site operator $B\ge0$ whose translated copies commute
    pairwise on the periodic $N$-site chain:
    \[
        [B_{i,i+1},B_{j,j+1}]=0.
    \]
\end{definition}

\begin{definition}[Single-bond commuting form]
    \label{def:mpdo_has_commuting_form}
    \lean{MPOTensor.HasCommutingForm}
    \leanok
    \uses{def:mpdo_commuting_form_data, def:mpo_tensor}
    An MPO tensor $M$ has the single-bond commuting-form property if for
    every $N \ge 2$ there is a two-site operator $B$ as in
    Definition~\ref{def:mpdo_commuting_form_data} and a constant $c>0$ such
    that
    \[
        \rho^{(N)}(M)=c\prod_{i=0}^{N-1}B_{i,i+1}.
    \]
\end{definition}

\begin{definition}[Single-bond product predicate]
    \label{def:mpdo_has_commuting_bond_product}
    \lean{MPOTensor.CommutingBondProductData}
    \lean{MPOTensor.HasCommutingBondProductAt}
    \lean{MPOTensor.HasCommutingBondProduct}
    \leanok
    \uses{def:mpdo_has_commuting_form}
    This condition asks, for every chain length $N\geq2$, for one positive
    bond on the full two-site space and one positive normalization such that
    \[
        \rho^{(N)}(M)=c\prod_{i=0}^{N-1}B_{i,i+1},
        \qquad c>0,\qquad B\ge0,
        \qquad [B_{i,i+1},B_{j,j+1}]=0.
    \]
    This auxiliary condition contains no explicit sector labels or
    multiplicities.
\end{definition}

\begin{theorem}[Single-bond predicate iff single-bond commuting form]
    \label{thm:mpdo_has_commuting_bond_product_iff_has_commuting_form}
    \lean{MPOTensor.hasCommutingBondProduct_iff_hasCommutingForm}
    \leanok
    \uses{def:mpdo_has_commuting_bond_product,
        def:mpdo_has_commuting_form}
    The single-bond product predicate is equivalent to the existence of
    single-bond commuting-form data on every finite chain. This is an
    equivalence between two presentations of the same condition.
\end{theorem}

\begin{proof}\leanok
    Expanding the two definitions gives the same identity
    $\rho^{(N)}(M)=c\prod_{i=0}^{N-1}B_{i,i+1}$ with $c>0$.
\end{proof}

\begin{definition}[Single-bond predicate with doubled-index idempotence]
    \label{def:mpdo_has_commuting_bond_product_zcl}
    \lean{MPOTensor.HasCommutingBondProductWithZCL}
    \leanok
    \uses{def:mpdo_has_commuting_bond_product, def:mpo_is_zcl}
    This is the conjunction of the single-bond condition with the
    doubled-index transfer identity
    \[
        \E_M\circ\E_M=\E_M.
    \]
    The source physical-trace condition is
    Definition~\ref{def:mpo_source_zcl}.
\end{definition}

\begin{theorem}[Restricted predicate with idempotence iff restricted form with it]
    \label{thm:mpdo_has_commuting_bond_product_zcl_iff_has_commuting_form_and_is_zcl}
    \lean{MPOTensor.hasCommutingBondProductWithZCL_iff_hasCommutingForm_and_isZCL}
    \leanok
    \uses{def:mpdo_has_commuting_bond_product_zcl,
        thm:mpdo_has_commuting_bond_product_iff_has_commuting_form}
    The single-bond condition with doubled-index transfer idempotence is
    equivalent to the conjunction of the single-bond commuting-form
    property with that same idempotence condition.
\end{theorem}

\begin{proof}\leanok
    Expand the restricted conjunction and apply
    Theorem~\ref{thm:mpdo_has_commuting_bond_product_iff_has_commuting_form}.
\end{proof}

\begin{theorem}[A single bond gives one GSNNCH sector]
    \label{thm:mpdo_commuting_bond_product_has_gsnnch_form}
    \lean{MPOTensor.GSNNCHData.ofCommutingFormData}
    \lean{MPOTensor.CommutingFormData.hasGSNNCHFormAt_of_realizes}
    \lean{MPOTensor.hasGSNNCHFormAt_of_hasCommutingBondProductAt}
    \lean{MPOTensor.hasGSNNCHForm_of_hasCommutingBondProduct}
    \leanok
    \uses{def:mpdo_source_gsnnch,
        def:mpdo_has_commuting_bond_product}
    A single positive commuting bond is the one-sector instance of the
    explicit GSNNCH decomposition, with sector projection equal to the
    identity and multiplicity one.
\end{theorem}

\begin{proof}\leanok
    Take one sector, its one-site projection to be the identity, its
    multiplicity to be one, and its bond to be the given bond.
\end{proof}
